%% 第1巻　かぞえる・はかる　本文（第5回編集委員会で完成）
%% 序章 p.10〜13、第1部 p.14〜85、第2部 p.86〜127、第3部 p.128〜179、第4部 p.180〜231、巻末 p.232〜252
%% 章は \bkinputchap{ch1-NN} で読む（試し組みのときに、まだない章を（仮）にできるように）
\providecommand\bkinputchap[1]{\input{#1}}
\begingroup\bkvolonebody
\clearpage\setcounter{page}{10}
%% 1-0 序章　かずを つかうと、なにが いいの？（4ページ、p.10〜13）
\makeatletter
\gdef\bk@partlabel{\R{第}{だい}1\R{巻}{かん}}\gdef\bk@partname{かぞえる・はかる}\gdef\bk@tabpos{0}
\makeatother
\bkchapter{1}{0}{かずを つかうと、\\なにが いいの？}{\bklevel{1}}{かずを つかうと、なにが いいの？}

\noindent あしたは えんそく。 ソラと ミナが、 おやつの あめを みせあって います。

\begin{bkfigure}
\begin{tikzpicture}[x=1mm, y=1mm]
  \fill[bkYellowBg, rounded corners=4mm] (0,0) rectangle (132,62);
  % ソラ
  \bkChar{sora}{(19,4)}{34mm}
  \node[rectangle callout, callout absolute pointer={(30,34)}, draw=bkTheme, line width=0.8pt, fill=white,
    rounded corners=2.5mm, inner xsep=2.5mm, inner ysep=1.6mm, font=\small, align=left, anchor=south west] at (4,42) {あめ、 たくさん\\もって きたよ！};
  \fill[white, rounded corners=2mm] (36,6) rectangle (62,30);
  \foreach \i in {0,...,7} {\pgfmathtruncatemacro\cx{mod(\i,3)}\pgfmathtruncatemacro\cy{div(\i,3)}
    \bkcandy{(41+\cx*8,25-\cy*7)}{1.6}{bkRoute4}}
  % ミナ
  \bkChar{mina}{(113,4)}{34mm}
  \node[rectangle callout, callout absolute pointer={(102,34)}, draw=bkTheme, line width=0.8pt, fill=white,
    rounded corners=2.5mm, inner xsep=2.5mm, inner ysep=1.6mm, font=\small, align=left, anchor=south east] at (128,42) {わたしも\\たくさん！};
  \fill[white, rounded corners=2mm] (66,6) rectangle (96,30);
  \foreach \i in {0,...,11} {\pgfmathtruncatemacro\cx{mod(\i,4)}\pgfmathtruncatemacro\cy{div(\i,4)}
    \bkcandy{(70.5+\cx*7,25-\cy*7)}{1.6}{bkRoute6}}
  \node[font=\sffamily\bfseries\Large, text=bkTheme] at (66,47) {？};
\end{tikzpicture}
\end{bkfigure}

\noindent ふたりとも 「たくさん」。 でも、 どっちが おおいのかな？\\
「たくさん」と いう ことばだけでは、 くらべられません。

%% ---------------- かずに すると、くらべられる（第11回：ペアと かぞえるの 2とおりで みせる） ----------------
\bksubsection{1}{かずに すると、 くらべられる}

\noindent あめを 1こずつ ならべて、 ソラの あめと ミナの あめを 1つずつ ペアに して みましょう。

\begin{bkfigure}
\begin{tikzpicture}[x=1mm, y=1mm]
  \fill[bkYellowBg, rounded corners=3mm] (-2,-8) rectangle (130,24);
  \node[anchor=west, font=\sffamily\bfseries\small, text=bkInk] at (0,15) {ソラ};
  \node[anchor=west, font=\sffamily\bfseries\small, text=bkInk] at (0,0) {ミナ};
  \foreach \i in {0,...,7} {\bkcandy{(24+\i*9,15)}{1.6}{bkRoute4}}
  \foreach \i in {0,...,11} {\bkcandy{(24+\i*9,0)}{1.6}{bkRoute6}}
  \foreach \i in {0,...,7} {\draw[bkGray, line width=0.8pt] (24+\i*9,12.2) -- (24+\i*9,2.8);}
  \draw[bkTheme, line width=1.2pt] (110.5,0) ellipse (18 and 4.8);
  \node[font=\sffamily\bfseries\small, text=bkTheme] at (110.5,15) {あまった 4こ};
  \node[font=\sffamily\footnotesize, text=bkGray!80!black] at (55.5,-5.6) {せんで むすんだ ペア};
\end{tikzpicture}
\end{bkfigure}

\noindent ミナの あめが 4こ あまりました。 あまった ほうが おおいね。

\noindent こんどは、 1こずつ ゆびで さしながら かぞえて みましょう。 ソラは \textbf{8こ}、 ミナは \textbf{12こ}。\\
かぞえる ときは、 1、 2、 3、 …と、 じゅんばんに いいますね。 8の あとに 9、 10、 11、 12と つづきます。 あとから でて くる かずの ほうが、 おおきい かずです。 だから、 12は 8より おおきい。 ならべなくても、 かずを きくだけで 「ミナの ほうが おおい」と わかります。

\begin{taisetsu}
かずに すると、 \textbf{くらべられる}。
\end{taisetsu}

\noindent かずの かぞえかたと くらべかたは、 1-1で ゆっくり やります。 12のような 10より おおきい かずは、 1-2で まなびます。

%% ---------------- p.11 つたえられる・さきが わかる ----------------
\bksubsection{2}{かずに すると、 \R{人}{ひと}に つたえられる}

\noindent\begin{minipage}[c]{0.62\linewidth}
「リンゴを たくさん かって きてね」と たのまれたら、 なんこ かえば いいのか こまって しまいます。 さっきの 「たくさん」も、 ソラは 8こ、 ミナは 12こでした。 「たくさん」は、 \R{人}{ひと}に よって ちがうのです。

「リンゴを 6こ かって きてね」なら、 だれが いっても、 おなじ かずだけ かって こられます。
\end{minipage}\hfill
\begin{minipage}[c]{0.35\linewidth}\centering
\begin{tikzpicture}[x=1mm, y=1mm]
  \fill[white, draw=bkLine, line width=0.8pt, rounded corners=3mm] (0,0) rectangle (46,34);
  \bkplate{(23,9)}{38}
  \foreach \x/\y in {10/12,18/14,26/14.5,34/13.5,14/20.5,23/21.5} {\bkapple{(\x,\y)}{3.6}{bkAppleRed}}
  \node[font=\sffamily\bfseries\small, text=bkTheme] at (37,27) {6こ};
\end{tikzpicture}
\end{minipage}

\bksubsection{3}{かずに すると、 この さきが わかる}

\noindent\begin{minipage}[c]{0.6\linewidth}
\R{水}{みず}そうに バケツで \R{水}{みず}を いれます。 バケツ 1ぱいで、 \R{水}{みず}は \R{線}{せん}（\R{目}{め}もり） 1つ\R{分}{ぶん} あがります。

いまは \R{目}{め}もり 6。 いっぱいは \R{目}{め}もり 10です。 あと なんばいで いっぱいに なるかな？

いまの 6の つぎから、 7、 8、 9、 10と \R{目}{め}もりを かぞえると 4つ。 バケツ 1ぱいで \R{目}{め}もり 1つ だから、 あと 4はいです。 \R{水}{みず}を いれる まえに、 かずで わかりました。
\end{minipage}\hfill
\begin{minipage}[c]{0.37\linewidth}\centering
\begin{tikzpicture}[x=1mm, y=1mm]
  % 水そう
  \fill[bkBlue!25] (2,0.6) rectangle (34,30);
  \draw[bkInk!70, line width=1.2pt, line join=round] (1.4,50) -- (1.4,0) -- (34.6,0) -- (34.6,50);
  \foreach \k in {1,...,10} {\draw[bkInk, line width=0.6pt] (34.6,\k*5) -- (30.6,\k*5);
    \node[anchor=west, font=\sffamily\bfseries\footnotesize, text=bkInk] at (35.4,\k*5) {\k};}
  \draw[bkRed, line width=1pt, dash pattern=on 2pt off 1.5pt] (1.4,50) -- (34.6,50);
  \node[anchor=south, font=\sffamily\bfseries\footnotesize, text=bkRed] at (18,50.5) {いっぱい};
  \foreach \k in {7,...,10} {\draw[bkOrange, line width=1.1pt, -{Stealth[length=3pt]}] (25,{\k*5-4.3}) -- (25,{\k*5-0.7});}
  \node[anchor=east, font=\sffamily\bfseries\footnotesize, text=bkOrange!80!black, align=right] at (23,42.5) {あと\\4はい};
  % バケツ
  \begin{scope}[shift={(-2,60)}, rotate=-35]
    \fill[bkGray!50] (0,0) -- (12,0) -- (10.5,-10) -- (1.5,-10) -- cycle;
    \draw[bkInk!70, line width=0.8pt] (0,0) .. controls (2,5) and (10,5) .. (12,0);
  \end{scope}
\end{tikzpicture}
\end{minipage}

\par\vspace{8pt}
\noindent カレンダーでも、 おなじように さきが わかります。 きょうは \R{5日}{いつか}。 えんそくは \R{8日}{ようか}です。

\begin{bkfigure}
\begin{tikzpicture}[x=1mm, y=1mm]
  \fill[bkTheme!6, rounded corners=3mm] (0,-9) rectangle (132,31);
  \foreach \d [count=\i from 0] in {3,4,5,6,7,8,9} {
    \filldraw[fill=white, draw=bkGray, line width=0.7pt] (8+\i*17,0) rectangle ++(15,15);
    \node[font=\sffamily\bfseries\large, text=bkInk] at (15.5+\i*17,7.5) {\d};}
  \draw[bkBlue, line width=1.3pt] (49.5,7.5) circle (6);
  \node[font=\sffamily\bfseries\footnotesize, text=bkBlue!80!black] at (49.5,-4.5) {きょう};
  \node[font=\sffamily\bfseries\footnotesize, text=bkTheme] at (100.5,-4.5) {えんそく};
  \draw[bkTheme, line width=1.3pt] (100.5,7.5) circle (6);
  \foreach \i in {0,1,2} {\draw[bkOrange, line width=1.3pt, -{Stealth[length=5pt]}] (51+\i*17,16.5) .. controls (55+\i*17,26) and (61+\i*17,26) .. (65+\i*17,16.5);
    \node[circle, fill=bkOrange, text=white, font=\sffamily\bfseries\footnotesize, minimum size=4.6mm, inner sep=0pt] at (58+\i*17,27) {\the\numexpr\i+1\relax};}
\end{tikzpicture}
\end{bkfigure}

\noindent 5の つぎから、 6、 7、 8と かぞえると 3つ。 あと 3かい ねると、 えんそくの \R{日}{ひ}です。 かずと じゅんばんが わかって いれば、 その \R{日}{ひ}が くる まえに わかります。

\clearpage
%% ---------------- きかいにも つたえられる（第11回で追加） ----------------
\bksubsection{4}{かずなら、 きかいにも つたえられる}

\noindent きかいは、 「たくさん」や 「すこし」が わかりません。 でも、 かずで つたえると、 そのとおりに うごきます。

\begin{bkfigure}
\begin{tikzpicture}[x=1mm, y=1mm,
  btn/.style={circle, draw=bkGray, line width=0.8pt, fill=white, minimum size=7.5mm, inner sep=0pt, font=\sffamily\bfseries\small}]
  % エレベーターの ボタン
  \fill[white, draw=bkLine, line width=0.8pt, rounded corners=3mm] (0,0) rectangle (62,58);
  \fill[bkGray!25, rounded corners=2mm] (16,10) rectangle (46,54);
  \fill[bkInk!85, rounded corners=1mm] (21,44) rectangle (41,51);
  \node[font=\sffamily\bfseries\small, text=bkYellow] at (31,47.5) {5かい};
  \foreach \n/\x/\y in {1/24/15,2/38/15,3/24/25,4/38/25,6/38/35} {\node[btn] at (\x,\y) {\n};}
  \node[btn, fill=bkYellow!60, draw=bkOrange] at (24,35) {5};
  \node[font=\sffamily\bfseries\small] at (31,4.5) {エレベーター};
  % でんしレンジ
  \begin{scope}[shift={(70,0)}]
  \fill[white, draw=bkLine, line width=0.8pt, rounded corners=3mm] (0,0) rectangle (62,58);
  \fill[bkGray!35, rounded corners=2mm] (5,14) rectangle (57,46);
  \fill[bkInk!80, rounded corners=1mm] (9,18) rectangle (40,42);
  \fill[bkGray!15, rounded corners=1mm] (43,18) rectangle (54,42);
  \fill[bkGreen!25] (44.5,34) rectangle (52.5,40);
  \node[font=\sffamily\bfseries\footnotesize, text=bkInk] at (48.5,37) {30};
  \foreach \y in {23,29} {\fill[bkGray!70] (48.5,\y) circle (1.8);}
  \node[font=\sffamily\bfseries\footnotesize, text=bkInk] at (31,52) {30びょう あたためる};
  \node[font=\sffamily\bfseries\small] at (31,4.5) {でんしレンジ};
  \end{scope}
\end{tikzpicture}
\end{bkfigure}

\begin{itemize}[label={\color{bkPurple}●}, leftmargin=1.4\zw, itemsep=0pt, topsep=0pt]
\item エレベーターに 「うえの ほう」と いっても、 どこで とまれば いいか わかりません。 「5」の ボタンを おすと、 5かいで とまります。
\item でんしレンジは、 「30」と かずで いれると、 30びょうだけ あたためます。
\end{itemize}

\noindent ロボットも ロケットも おなじです。 \R{人}{ひと}が かずで つたえるから、 きまった とおりに うごけるのです。

\begin{taisetsu}
かずに すると、 \textbf{くらべられる}。 \textbf{つたえられる}。\\
\textbf{この さきが わかる}。 \textbf{きかいも うごかせる}。
\end{taisetsu}


\clearpage
%% ---------------- p.12 この本で できるように なる こと ----------------
\bksection{goal}{この \R{本}{ほん}で できるように なる こと}

\noindent この \R{本}{ほん}は、 4つの 「ぶ」に わかれて います。

\begin{bkfigure}
\begin{tikzpicture}[x=1mm, y=1mm,
  pb/.style={draw=bkVol, line width=1.2pt, fill=white, rounded corners=3mm, minimum width=63mm, minimum height=50mm}]
  \foreach \i/\x/\y in {1/0/53,2/69/53,3/0/0,4/69/0} {\node[pb, anchor=south west] (p\i) at (\x,\y) {};}
  \foreach \i/\x/\y/\t/\c/\pg in {
      1/0/53/{かずの はじまり}/{かぞえる・ くらべる・ たす・ ひく}/{1-1〜1-7},
      2/69/53/はかる/{ながさ・ かさ・ おもさ・ \R{時間}{じ|かん}を はかる}/{1-8〜1-11},
      3/0/0/{かけ\R{算}{ざん}と わり\R{算}{ざん}}/{\R{同}{おな}じ かずの まとまりで かぞえる・ \R{分}{わ}ける}/{1-12〜1-16},
      4/69/0/{\R{表}{ひょう}・グラフ・かたち}/{\R{表}{ひょう}や グラフに する・ かたちと \R{広}{ひろ}さ}/{1-17〜1-21}} {
    \fill[bkVol, rounded corners=2mm] (\x+3,\y+40) rectangle (\x+17,\y+47);
    \node[text=bkOn, font=\sffamily\bfseries\footnotesize] at (\x+10,\y+43.2) {\R{第}{だい}\i\R{部}{ぶ}};
    \node[anchor=west, font=\sffamily\bfseries\small, text=bkDeep] at (\x+19,\y+43.2) {\t};
    \node[anchor=north west, font=\footnotesize, text=bkInk, text width=57mm, align=left, inner sep=0pt] at (\x+4,\y+36.5) {\c};
    \node[anchor=south east, font=\sffamily\bfseries\footnotesize, text=bkGray] at (\x+61,\y+1.5) {\pg};}% 第11回：ページ数が かわるので、ページ番号ではなく 章の はんいを かく
  % 第1部：リンゴと 10のわく
  \foreach \x in {7,14,21} {\bkapple{(\x,65)}{3}{bkAppleRed}}
  \node[font=\sffamily\bfseries\small, text=bkTheme] at (28,65) {＋};
  \foreach \x in {35,42} {\bkapple{(\x,65)}{3}{bkAppleGreen}}
  \bktenframe{(5,54.5)}{3.2}{3}{bkAppleRed}{2}{bkAppleGreen}
  % 第2部：ものさし・時計・ます
  \begin{scope}[shift={(69,53)}]
    \fill[bkYellowBg, draw=bkYellow!80!black, line width=0.6pt, rounded corners=0.8mm] (4,14) rectangle (30,20);
    \foreach \k in {0,...,10} \draw[bkInk, line width=0.4pt] (5.5+\k*2.3,20) -- (5.5+\k*2.3,{\k==5 ? 17 : 18.4});
    \bkclock*{(38,15)}{7}{3}{0}
    \bkmasu{(10,1.5)}{7}{0.6}{0}
  \end{scope}
  % 第3部：おはじきの ならび
  \foreach \i in {0,...,3} \foreach \j in {0,1,2} {\bkohajiki{(8+\i*7,{6+\j*6})}{2.3}{bkRoute4}}
  \node[font=\sffamily\bfseries\small, text=bkInk] at (42,12) {3 × 4};
  % 第4部：ぼうグラフと まる
  \begin{scope}[shift={(69,0)}]
    \draw[bkInk!70, line width=0.6pt] (5,4) -- (34,4) (5,4) -- (5,22);
    \foreach \h/\x/\c in {10/9/bkRoute2,16/17/bkRoute6,7/25/bkRoute3} {\fill[\c] (\x,4) rectangle ++(5,\h);}
    \draw[bkRoute4, line width=1.2pt] (44,15) circle (7);
    \draw[bkRoute4, line width=0.8pt] (44,15) -- ++(30:7);
    \fill[bkRoute4] (44,15) circle (0.8);
  \end{scope}
\end{tikzpicture}
\end{bkfigure}

\noindent えの なかの 「＋」や 「×」は、 これから この \R{本}{ほん}で まなぶ しるしだよ。 いまは わからなくても だいじょうぶ。

\noindent どの \R{章}{しょう}も 「はてな？」から はじまります。 さいごには、 その ふしぎに かずで こたえられるように なるよ。

\noindent \R{章}{しょう}の おわりの 「ふりかえり」で、 できるように なった ことに チェックを つけよう。 \R{本}{ほん}の うしろの 「チェックひょう」にも、 まとめて あります。

\clearpage
%% ---------------- p.13 3人の かいわ ----------------
\bksection{reidai}{3\R{人}{にん}の なかまと いっしょに}

\bkfukidashi{sora}{かずって、 かぞえる ためだけの もの だと おもって いたよ。}
\bkfukidashi{mina}{はかる ときも、 くらべる ときも、 かずを つかうんだね。}
\bkfukidashi{delta}{そうだよ。 ロケットも ロボットも、 かずで うごいて いるんだ。 1から いっしょに はじめよう！}

\begin{bkfigure}
\begin{tikzpicture}[x=1mm, y=1mm]
  \fill[bkTheme!6, rounded corners=4mm] (0,-1) rectangle (132,40);
  % 1巻から9巻への道
  \draw[bkTheme, line width=2.4pt, line cap=round] (10,12) .. controls (40,12) and (50,30) .. (80,28) .. controls (100,27) and (110,20) .. (122,26);
  \foreach \x/\y/\k in {10/12/1,27/13.5/2,42/19/3,57/25.5/4,72/28.3/5,87/27.6/6,101/25.3/7,113/22.6/8,122/26/9} {
    \fill[white, draw=bkVol\k, line width=1.4pt] (\x,\y) circle (2.8);
    \node[font=\sffamily\bfseries\scriptsize, text=bkVol\k!85!black] at (\x,\y) {\k};}
  \node[fill=bkTheme, text=white, rounded corners=1mm, font=\sffamily\bfseries\footnotesize, inner sep=1.5pt] at (10,20) {いま ここ};
  \node[anchor=north west, font=\footnotesize, text=bkInk, align=left] at (4,8) {1\R{巻}{かん}：かぞえる・ はかる};
  \node[anchor=north east, font=\footnotesize, text=bkInk, align=right] at (130,18) {9\R{巻}{かん}：ロボットや\\ロケットを うごかす すうがく};
\end{tikzpicture}
\end{bkfigure}

\noindent この \R{本}{ほん}は、 9さつ つづく \R{本}{ほん}の 1さつめです。 1さつずつ すすむと、 さいごには、 ロボットや ロケットを うごかす すうがくまで とどきます。

\noindent この \R{本}{ほん}の つかいかたは、 8ページを みてね。 わからない ところが あっても だいじょうぶ。 ゆっくり、 いっしょに すすもう。


%% ======================== 第1部 かずの はじまり ========================
\bkPartOpener{1}{\R{第}{だい}1\R{巻}{かん}\hspace{0.8em}\R{第}{だい}1\R{部}{ぶ}}{かずの はじまり}%
  {172/{おさらと リンゴ、 どっちが おおい？},
   160/{10より おおきな かずは、 どう かぞえる？},
   148/{なにも ないのに、 「0」と かくのは なぜ？}}%
  {% 10のわくと ものさし（白い線）
   \begin{scope}[white, line width=1.1pt]
     \foreach \i in {0,...,9} {\pgfmathtruncatemacro\cx{mod(\i,5)}\pgfmathtruncatemacro\cy{1-div(\i,5)}
       \draw (124+\cx*8,122+\cy*8) rectangle ++(8,8);
       \ifnum\i<7 \fill (128+\cx*8,126+\cy*8) circle (2.6);\fi}
     \draw[rounded corners=1mm] (124,107) rectangle (168,116);
     \foreach \x in {127,131,...,165} \draw[line width=0.7pt] (\x,116) -- (\x,113.5);
   \end{scope}}%
  {かずを つかうと、 くらべたり、\\ つたえたり できるよ。 この ぶでは、\\ かずの きほんと、 たし\R{算}{ざん}・ひき\R{算}{ざん}を まなぶよ。}
\bkPartContents{1}{%% 自動生成（gen.py）
\bktocchap{1-1}{いくつ あるかな？}{おさらと リンゴ、どっちが おおい？}{2}{1}{1}{18}{C}
\bktocchap{1-2}{10の まとまりを つくろう}{10ずつ たばねると、\R{大}{おお}きな かずも すぐ わかる}{3}{1}{1}{34}{C}
\bktocchap{1-3}{なにも ないときの「0」}{ものさしの はじまりは、どうして 0なの？}{6}{1}{1}{52}{C}
\bktocchap{1-4}{たし\R{算}{ざん}は「あわせる」「ふえる」}{リンゴ 3こと 2こ、あわせて いくつ？}{2}{1}{1}{66}{C}
\bktocchap{1-5}{ひき\R{算}{ざん}は「のこり」「ちがい」「かわった \R{量}{りょう}」}{きのうより、どれだけ のびた？}{1}{1}{1}{80}{C}
\bktocchap{1-6}{かずの \R{線}{せん}：\R{数直線}{すう|ちょく|せん}}{ものさしの \R{上}{うえ}を、すすんだり もどったり}{6}{1}{1}{94}{C}
\bktocchap{1-7}{\R{大}{おお}きな かずの たし\R{算}{ざん}・ひき\R{算}{ざん}}{36 ＋ 27 は？ 10の まとまりで かんがえよう}{2}{1}{1}{112}{C}
}%
{「たくさん」より 「12こ」の ほうが、 よく つたわります。 この ぶで まなぶ 「10の まとまり」と 「たし\R{算}{ざん}・ひき\R{算}{ざん}」は、 さいごの \R{本}{ほん}まで ずっと つかいます。}
\bkpartbegin{1}{かずの はじまり}
\bkinputchap{ch1-01}
\bkinputchap{ch1-02}
\bkinputchap{ch1-03}
\bkinputchap{ch1-04}
\bkinputchap{ch1-05}
\bkinputchap{ch1-06}
\bkinputchap{ch1-07}

%% ======================== 第2部 はかる ========================
\bkPartOpener{2}{\R{第}{だい}1\R{巻}{かん}\hspace{0.8em}\R{第}{だい}2\R{部}{ぶ}}{はかる}%
  {172/{えんぴつ \R{何本}{なん|ぼん}\R{分}{ぶん}？ \R{人}{ひと}に よって ちがったら？},
   160/{\R{大}{おお}きい ほうが、 おもいとは かぎらない？},
   148/{「3\R{時}{じ}」と 「3\R{時間}{じ|かん}」は、 どこが ちがう？}}%
  {% ものさしと 時計と ます（白い線）
   \begin{scope}[white, line width=1.1pt]
     \draw[rounded corners=1mm] (118,107) rectangle (170,116);
     \foreach \k in {0,...,20} {\pgfmathsetmacro\bkTL{ifthenelse(mod(\k,10)==0,4.5,ifthenelse(mod(\k,5)==0,3.4,2.2))}
       \draw[line width=0.7pt] (121+\k*2.3,116) -- ++(0,-\bkTL);}
     \draw (136,138) circle (12);
     \foreach \h in {1,...,12} \draw[line width=0.8pt] ({136+10.5*cos(90-30*\h)},{138+10.5*sin(90-30*\h)}) -- ({136+12*cos(90-30*\h)},{138+12*sin(90-30*\h)});
     \draw[line width=1.6pt, line cap=round] (136,138) -- (142,138);
     \draw[line width=1.1pt, line cap=round] (136,138) -- (136,147.5);
     \draw[line join=round] (156,148) -- (156,124) -- (172,124) -- (172,148);
     \fill[white, opacity=0.35] (156.6,124.6) rectangle (171.4,138);
     \foreach \k in {1,...,4} \draw[line width=0.7pt] (172,124+\k*4.8) -- ++(-2.4,0);
   \end{scope}}%
  {ながさ・ かさ・ おもさ・ \R{時間}{じ|かん}を、 みんなが\\ おなじ \R{目}{め}もりで はかれるように なるよ。\\ 「はかる」は、 \R{工学}{こう|がく}の いちばん はじめの \R{道具}{どう|ぐ}だよ。}
\bkPartContents{2}{%% 自動生成（gen.py）
\bktocchap{1-8}{ながさを はかろう}{えんぴつ \R{何本}{なん|ぼん}\R{分}{ぶん}？ \R{人}{ひと}に よって ちがったら こまるね}{7}{1}{1}{132}{C}
\bktocchap{1-9}{かさと おもさ}{\R{大}{おお}きい ほうが、おもいとは かぎらない？}{7}{1}{1}{150}{C}
\bktocchap{1-10}{\R{時}{じ}こくと \R{時間}{じ|かん}}{「3\R{時}{じ}」と「3\R{時間}{じ|かん}」は、どこが ちがう？}{1}{1}{1}{166}{C}
\bktocchap{1-11}{\R{目}{め}もりを よもう}{\R{目}{め}もりと \R{目}{め}もりの あいだは、どう よむ？}{7}{1}{1}{184}{C}
}%
{だれが はかっても おなじ かずに なるように、 cm・L・g・\R{分}{ふん}などの 「たんい」が きまって います。 この ぶで、 1つずつ まなびます。 \R{目}{め}もりを よむ \R{力}{ちから}は、 この さきの \R{本}{ほん}で、 じっけんや グラフを よむ ときに ずっと つかいます。}
\bkpartbegin{2}{はかる}
\bkinputchap{ch1-08}
\bkinputchap{ch1-09}
\bkinputchap{ch1-10}
\bkinputchap{ch1-11}

%% ======================== 第3部 かけ算と わり算 ========================
\bkPartOpener{3}{\R{第}{だい}1\R{巻}{かん}\hspace{0.8em}\R{第}{だい}3\R{部}{ぶ}}{かけ\R{算}{ざん}と わり\R{算}{ざん}}%
  {172/{\R{同}{おな}じ かずずつ ならんだ ものを、 はやく かぞえるには？},
   160/{\R{九九}{く|く}の \R{表}{ひょう}には、 どんな きまりが ある？},
   148/{12こを 3\R{人}{にん}で \R{分}{わ}けると、 ひとり いくつ？}}%
  {% おはじきの ならび（白い丸）
   \begin{scope}[white]
     \foreach \i in {0,...,5} \foreach \j in {0,...,3} {\fill (124+\i*8,112+\j*8) circle (2.7);}
     \draw[line width=1.1pt, rounded corners=2mm] (119,107) rectangle (169,143);
     \draw[line width=1.1pt, dash pattern=on 3pt off 2pt] (127.5,107) -- (127.5,143) (119,115.5) -- (169,115.5);
   \end{scope}}%
  {\R{同}{おな}じ かずの まとまりを つかうと、\\ たくさんの ものも はやく かぞえられるよ。\\ \R{分}{わ}ける ときも、 \R{九九}{く|く}が \R{役}{やく}に たつんだ。}
\bkPartContents{3}{%% 自動生成（gen.py）
\bktocchap{1-12}{かけ\R{算}{ざん}は「\R{同}{おな}じ かずの まとまり」}{おはじきを ならべると、\R{四角}{し|かく}が できるよ}{2}{1}{1}{202}{C}
\bktocchap{1-13}{\R{九九}{く|く}の \R{表}{ひょう}を たんけんしよう}{\R{表}{ひょう}を ななめに \R{見}{み}ると、ふしぎな きまりが…}{5}{1}{1}{218}{C}
\bktocchap{1-14}{わり\R{算}{ざん}は「\R{同}{おな}じ かずずつ \R{分}{わ}ける」}{12こを 3\R{人}{にん}で \R{分}{わ}けると、ひとり いくつ？}{1}{1}{1}{234}{C}
\bktocchap{1-15}{あまりの ある わり\R{算}{ざん}}{100\R{日}{にち}\R{後}{ご}は、\R{何曜日}{なん|よう|び}？}{4}{1}{1}{250}{C}
\bktocchap{1-16}{1より \R{小}{ちい}さい かず：\R{小数}{しょう|すう}と \R{分数}{ぶん|すう}}{ピザを \R{半分}{はん|ぶん}の \R{半分}{はん|ぶん}に すると？}{7}{1}{1}{266}{C}
}%
{この ぶで まなぶ 「1つ\R{分}{ぶん} × いくつ\R{分}{ぶん}」と いう かんがえかたは、 この さきの \R{本}{ほん}で 「はやさ」や 「\R{面積}{めん|せき}」を かんがえる もとに なります。 \R{九九}{く|く}の \R{表}{ひょう}は、 7\R{巻}{かん}の 「\R{行列}{ぎょう|れつ}」の たねです。}
\bkpartbegin{3}{かけ\R{算}{ざん}と わり\R{算}{ざん}}
\bkinputchap{ch1-12}
\bkinputchap{ch1-13}
\bkinputchap{ch1-14}
\bkinputchap{ch1-15}
\bkinputchap{ch1-16}

%% ======================== 第4部 表・グラフ・かたち ========================
\bkPartOpener{4}{\R{第}{だい}1\R{巻}{かん}\hspace{0.8em}\R{第}{だい}4\R{部}{ぶ}}{\R{表}{ひょう}・グラフ・かたち}%
  {172/{いちばん \R{人気}{にん|き}の きゅう\R{食}{しょく}は、 どれ？},
   160/{\R{気温}{き|おん}が いちばん はやく \R{上}{あ}がったのは \R{何時}{なん|じ}ごろ？},
   148/{タイヤが \R{四角}{し|かく}だったら、 どうなる？}}%
  {% ぼうグラフと おれ線と 円（白い線）
   \begin{scope}[white, line width=1.1pt]
     \draw (118,108) -- (150,108) (118,108) -- (118,142);
     \foreach \x/\h in {122/14,130/24,138/18} {\fill[opacity=0.85] (\x,108) rectangle ++(5,\h);}
     \draw[line width=1.3pt] (120,124) -- (128,130) -- (136,140) -- (144,136) -- (150,128);
     \foreach \p in {(120,124),(128,130),(136,140),(144,136),(150,128)} \fill \p circle (1.3);
     \draw (162,124) circle (10);
     \draw (162,124) -- ++(35:10);
     \fill (162,124) circle (1);
   \end{scope}}%
  {しらべた ことを \R{表}{ひょう}や グラフに すると、\\ ひとめで わかるように なるよ。 かたちや\\ \R{広}{ひろ}さも、 かずで あらわせるんだ。}
\bkPartContents{4}{%% 自動生成（gen.py）
\bktocchap{1-17}{\R{表}{ひょう}と ぼうグラフ}{いちばん \R{人気}{にん|き}の きゅう\R{食}{しょく}は？ \R{一目}{ひと|め}で わかる}{5}{1}{1}{284}{C}
\bktocchap{1-18}{おれ\R{線}{せん}グラフ：かわり\R{方}{かた}を \R{見}{み}る}{\R{気温}{き|おん}が いちばん はやく \R{上}{あ}がったのは \R{何時}{なん|じ}ごろ？}{1}{1}{1}{304}{C}
\bktocchap{1-19}{かたち：\R{三角}{さん|かく}・\R{四角}{し|かく}・\R{円}{えん}}{タイヤが \R{四角}{し|かく}だったら、どうなる？}{4}{1}{1}{322}{C}
\bktocchap{1-20}{\R{角}{かく}の \R{大}{おお}きさ}{1\R{回転}{かい|てん}は、どれだけ まわった ことに なる？}{4}{1}{1}{340}{C}
\bktocchap{1-21}{\R{広}{ひろ}さを くらべよう：\R{面積}{めん|せき}}{どっちの へやが \R{広}{ひろ}い？ ならべて くらべよう}{2}{1}{1}{358}{C}
}%
{グラフの 「かたむき」は、 ものが かわる はやさです。 5\R{巻}{かん}の 「\R{微分}{び|ぶん}」に つながります。 \R{角}{かく}の \R{大}{おお}きさと \R{面積}{めん|せき}は、 ロボットの うごきや、 \R{電気}{でん|き}・\R{水}{みず}の ながれを かんがえる もとに なります。}
\bkpartbegin{4}{\R{表}{ひょう}・グラフ・かたち}
\bkinputchap{ch1-17}
\bkinputchap{ch1-18}
\bkinputchap{ch1-19}
\bkinputchap{ch1-20}
\bkinputchap{ch1-21}

%% ======================== 巻末 ========================
%% 第1巻 巻末（こたえ・きごうカード・たんいのひょう・チェックひょう・さくいん・奥付）
\makeatletter
%% ---------------- こたえ ----------------
\bkAnswerPages

%% ---------------- きごうカード ----------------
\bkKigoPagesOne

%% ---------------- たんいのひょう ----------------
\bkbacktitle{たんいのひょう}
{\small この \R{本}{ほん}で でてきた たんいと、 その \R{大}{おお}きさの めやすです。}\par\vspace{2mm}

\def\bkuhead#1#2{\par\addvspace{3mm}\noindent{\sffamily\bfseries\large\color{bkTheme}#1}\hspace{0.6\zw}{\small\color{bkGray}#2}\par\nobreak\vspace{1mm}}
\newcommand\bkutable[1]{\begin{bkcenter}\small
\begin{tblr}{colspec={Q[c,m,wd=22mm] Q[l,m,wd=46mm] Q[l,m,wd=58mm]}, row{1}={bg=bkVol!14, font=\sffamily\bfseries\footnotesize},
  hlines={0.4pt,bkLine}, rowsep=2pt}
たんい & よみかた & めやす・ くらべかた \\
#1
\end{tblr}\end{bkcenter}}

\bkuhead{ながさ}{（1-8）}
\bkutable{%
1\,mm & ミリメートル & 1\,cm を 10に \R{分}{わ}けた 1つ \\
1\,cm & センチメートル & ゆびの はばほど。 1\,cm ＝ 10\,mm \\
1\,m & メートル & りょう\R{手}{て}を ひろげた ほど。 1\,m ＝ 100\,cm \\}

\bkuhead{かさ}{（1-9）}
\bkutable{%
1\,mL & ミリリットル & \R{小}{こ}さじ 1ぱいは 5\,mL \\
1\,dL & デシリットル & コップ \R{半分}{はん|ぶん}ほど。 1\,dL ＝ 100\,mL \\
1\,L & リットル & \R{牛乳}{ぎゅう|にゅう}パック 1\R{本}{ほん}。 1\,L ＝ 10\,dL ＝ 1000\,mL \\}

\bkuhead{おもさ}{（1-9）}
\bkutable{%
1\,g & グラム & 1\R{円}{えん}\R{玉}{だま} 1まい \\
1\,kg & キログラム & \R{水}{みず} 1\,L。 1\,kg ＝ 1000\,g \\}

\clearpage
\bkuhead{\R{時間}{じ|かん}}{（1-10）}
\bkutable{%
1\R{秒}{びょう} & びょう & 「いち」と いう あいだの \R{時間}{じ|かん} \\
1\R{分}{ぷん} & ふん・ぷん & 1\R{分}{ぷん} ＝ 60\R{秒}{びょう} \\
1\R{時間}{じ|かん} & じかん & 1\R{時間}{じ|かん} ＝ 60\R{分}{ぷん} \\
1\R{日}{にち} & にち & 1\R{日}{にち} ＝ 24\R{時間}{じ|かん} \\}

\bkuhead{あたたかさ（\R{温度}{おん|ど}）}{（1-11）}
\bkutable{%
1\,℃ & ど & \R{水}{みず}が こおるのは 0\,℃、 にえるのは 100\,℃ \\}

\bkuhead{\R{角}{かく}の \R{大}{おお}きさ}{（1-20）}
\bkutable{%
1° & ど & 1\R{回転}{かい|てん}を 360に \R{分}{わ}けた 1つ。 \R{直角}{ちょっ|かく} ＝ 90°、 1\R{回転}{かい|てん} ＝ 360° \\}

\bkuhead{\R{広}{ひろ}さ（\R{面積}{めん|せき}）}{（1-21）}
\bkutable{%
1\,cm² & へいほう センチメートル & 1\,cm の \R{正方形}{せい|ほう|けい}の \R{広}{ひろ}さ \\
1\,m² & へいほう メートル & 1\,m の \R{正方形}{せい|ほう|けい}の \R{広}{ひろ}さ \\}

\clearpage
\bkuhead{\R{九九}{く|く}の \R{表}{ひょう}}{（1-13）}
\begin{bkcenter}
\begin{tikzpicture}[x=1mm, y=1mm]
  \foreach \a in {1,...,9} {
    \fill[bkVol!18] (\a*13,-1*11) rectangle ++(13,11);
    \node[font=\sffamily\bfseries\small, text=bkDeep] at (\a*13+6.5,-5.5) {\a};
    \fill[bkVol!18] (0,-\a*11-11) rectangle ++(13,11);
    \node[font=\sffamily\bfseries\small, text=bkDeep] at (6.5,-\a*11-5.5) {\a};
    \foreach \b in {1,...,9} {\pgfmathtruncatemacro\pr{\a*\b}
      \ifnum\a=\b \fill[bkYellow!25] (\b*13,-\a*11-11) rectangle ++(13,11);\fi
      \node[font=\small, text=bkInk] at (\b*13+6.5,-\a*11-5.5) {\pr};}}
  \fill[bkVol!35] (0,-11) rectangle ++(13,11);
  \node[font=\sffamily\bfseries\footnotesize, text=bkDeep] at (6.5,-5.5) {×};
  \draw[bkLine, line width=0.4pt] (0,0) grid[xstep=13, ystep=11] (130,-110);
  \draw[bkVol, line width=1pt] (0,0) rectangle (130,-110) (13,0) -- (13,-110) (0,-11) -- (130,-11);
\end{tikzpicture}
\end{bkcenter}
{\small \R{左}{ひだり}の たての \R{数}{かず}が 「1つ\R{分}{ぶん}」、 \R{上}{うえ}の よこの \R{数}{かず}が 「いくつ\R{分}{ぶん}」です。 きいろの ななめの ますは、 \R{同}{おな}じ \R{数}{かず}どうしの かけ\R{算}{ざん}です。}

\vspace{4mm}
\bkuhead{\R{九九}{く|く}の きまり}{（1-13）}
\begin{bkcenter}
\begin{tikzpicture}[x=1mm, y=1mm]
  \foreach \y/\c/\t/\b/\e in {
    0/bkRoute2/{ふえかたの\\きまり}/{いくつ\R{分}{ぶん}が 1 ふえると、 1つ\R{分}{ぶん}だけ ふえる。}/{3 × 4 ＝ 12\enspace → \enspace 3 × 5 ＝ 15},
    -17/bkRoute4/{\R{入}{い}れかえても\\\R{同}{おな}じ}/{1つ\R{分}{ぶん}と いくつ\R{分}{ぶん}を \R{入}{い}れかえても \R{同}{おな}じ。}/{3 × 4 ＝ 4 × 3},
    -34/bkRoute5/{ななめの\\きまり}/{\R{同}{おな}じ かずどうしの かけ\R{算}{ざん}が ならぶ。}/{1 × 1、\enspace 2 × 2、\enspace 3 × 3、\enspace …}} {
    \fill[\c!12, rounded corners=2mm] (0,\y) rectangle ++(130,-15);
    \fill[\c, rounded corners=2mm] (0,\y) rectangle ++(3,-15);
    \node[anchor=west, font=\sffamily\bfseries\small, text=\c!80!black, inner sep=0pt, align=left] at (7,\y-7.5) {\t};
    \node[anchor=west, font=\small, text=bkInk, text width=93mm, inner sep=0pt, align=left] at (35,\y-7.5) {\b\\\sffamily\bfseries \e};}
\end{tikzpicture}
\end{bkcenter}

%% ---- 第5回で追加：かずの くらいと ひっ算（1-2・1-7） ----
\clearpage
\bkuhead{かずの くらい}{（1-2・1-7）}
\begin{bkcenter}
\begin{tikzpicture}[x=1mm, y=1mm]
  \foreach \i/\n/\c in {0/{\R{一万}{いち|まん}}/bkRoute5,1/{\R{千}{せん}}/bkGreen,2/{\R{百}{ひゃく}}/bkRoute4,3/{\R{十}{じゅう}}/bkBlue,4/{\R{一}{いち}}/bkOrange} {
    \fill[\c!14, rounded corners=1.5mm] (\i*25,0) rectangle ++(23,-30);
    \node[font=\sffamily\bfseries\small, text=\c!80!black] at (\i*25+11.5,-6) {\n};}
  \node[font=\fontsize{20}{20}\selectfont] at (11.5,-19) {10000};
  \node[font=\fontsize{20}{20}\selectfont] at (36.5,-19) {1000};
  \node[font=\fontsize{20}{20}\selectfont] at (61.5,-19) {100};
  \node[font=\fontsize{20}{20}\selectfont] at (86.5,-19) {10};
  \node[font=\fontsize{20}{20}\selectfont] at (111.5,-19) {1};
  \foreach \i in {0,...,3} {\draw[bk arrow, bkGray, line width=0.9pt] (\i*25+46,-33) to[bend left=35] node[below, font=\footnotesize, text=bkInk, inner sep=1.5pt] {10こで} (\i*25+21,-33);}
\end{tikzpicture}
\end{bkcenter}
{\small \R{上}{うえ}の \R{字}{じ}は くらいの \R{名前}{な|まえ}です（\R{一}{いち}の くらい、 \R{十}{じゅう}の くらい …）。 1が 10こで 10、 10が 10こで 100。 1つ \R{左}{ひだり}の くらいは、 10こ\R{分}{ぶん}です。}\par

\bkuhead{ひっ\R{算}{ざん}の しかた}{（1-7）}
\begin{bkcenter}
\begin{tikzpicture}[x=1mm, y=1mm, dg/.style={font=\fontsize{22}{22}\selectfont, text=bkInk, inner sep=0pt, anchor=base}]
  % たし算 36 ＋ 27
  \begin{scope}[shift={(19,0)}]
    \fill[bkBlue!10] (8,-31) rectangle ++(11,33); \fill[bkOrange!10] (19,-31) rectangle ++(11,33);
    \node[font=\sffamily\bfseries\scriptsize, text=bkBlue!80!black] at (13.5,4.5) {\R{十}{じゅう}};
    \node[font=\sffamily\bfseries\scriptsize, text=bkOrange!85!black] at (24.5,4.5) {\R{一}{いち}};
    \node[dg] at (13.5,-8) {3}; \node[dg] at (24.5,-8) {6};
    \node[dg] at (3,-19) {＋}; \node[dg] at (13.5,-19) {2}; \node[dg] at (24.5,-19) {7};
    \draw[bkInk, line width=1.1pt] (0,-22) -- (30,-22);
    \node[dg] at (13.5,-30) {6}; \node[dg] at (24.5,-30) {3};
    \node[font=\fontsize{11}{11}\selectfont, text=bkRed, inner sep=0pt] at (9.6,-12.8) {1};
  \end{scope}
  \node[anchor=north west, text width=64mm, font=\footnotesize, align=left] at (0,-36) {① \R{一}{いち}の くらい：6 ＋ 7 ＝ 13\\ ② 10を \R{十}{じゅう}の くらいへ（くり\R{上}{あ}がり）\\ ③ \R{十}{じゅう}の くらい：1 ＋ 3 ＋ 2 ＝ 6};
  % ひき算 52 － 27
  \begin{scope}[shift={(91,0)}]
    \fill[bkBlue!10] (8,-31) rectangle ++(11,33); \fill[bkOrange!10] (19,-31) rectangle ++(11,33);
    \node[font=\sffamily\bfseries\scriptsize, text=bkBlue!80!black] at (13.5,4.5) {\R{十}{じゅう}};
    \node[font=\sffamily\bfseries\scriptsize, text=bkOrange!85!black] at (24.5,4.5) {\R{一}{いち}};
    \node[dg] at (13.5,-8) {5}; \node[dg] at (24.5,-8) {2};
    \draw[bkRed, line width=0.9pt] (10.8,-9) -- (16.2,-1.2);
    \node[font=\fontsize{11}{11}\selectfont, text=bkRed, inner sep=0pt] at (17.6,-0.6) {4};
    \node[dg] at (3,-19) {－}; \node[dg] at (13.5,-19) {2}; \node[dg] at (24.5,-19) {7};
    \draw[bkInk, line width=1.1pt] (0,-22) -- (30,-22);
    \node[dg] at (13.5,-30) {2}; \node[dg] at (24.5,-30) {5};
  \end{scope}
  \node[anchor=north west, text width=64mm, font=\footnotesize, align=left] at (72,-36) {① 2から 7は ひけない\\ ② \R{十}{じゅう}の くらいから 10を 1つ もらう（くり\R{下}{さ}がり）。 12 － 7 ＝ 5\\ ③ \R{十}{じゅう}の くらい：4 － 2 ＝ 2};
\end{tikzpicture}
\end{bkcenter}
{\small くらいを たてに そろえて、 \R{一}{いち}の くらいから けいさんします。 こたえは、 たし\R{算}{ざん}と ひき\R{算}{ざん}で たしかめよう（36 ＋ 27 ＝ 63 なら、 63 － 27 ＝ 36）。}

%% ---------------- チェックひょう ----------------
\bkbacktitle{チェックひょう}
{\small \R{章}{しょう}の おわりの 「ふりかえり」を あつめた \R{表}{ひょう}です。 できるように なったら □に チェックを つけよう。 できなかったら、 その ページに もどって よみなおそう。}\par\vspace{2mm}
%% 自動生成（tools/mkchecklist.py）。各章の「ふりかえり」の項目を集めたもの
\begingroup\footnotesize
\def\bkcb{\tikz[baseline=-0.5ex]\draw[bkTheme, line width=0.7pt, rounded corners=0.5pt] (0,-0.3em) rectangle (0.8em,0.5em);}
\begin{longtblr}{colspec={Q[c,m,wd=9mm] X[l,m] Q[c,m,wd=10mm] Q[c,m,wd=9mm] Q[c,m,wd=9mm]}, width=\linewidth, rowhead=1, colsep=3pt,
  row{1}={bg=bkVol!14, font=\sffamily\bfseries\scriptsize}, hlines={0.4pt,bkLine}, rowsep=0.7pt}
\R{章}{しょう} & できるように なった こと & ページ & 1\R{回}{かい}め & 2\R{回}{かい}め \\
{\sffamily\bfseries\color{bkTheme}1-1} & 1つずつ ペアに して くらべられる。 & \pageref{bk:1-1:pea} & \bkcb & \bkcb \\
 & かぞえて、 すうじで かける。 & \pageref{bk:1-1:kazoeru} & \bkcb & \bkcb \\
 & おなじ かずは、 おなじ すうじ。 & \pageref{bk:1-1:onaji} & \bkcb & \bkcb \\
{\sffamily\bfseries\color{bkTheme}1-2} & 10ずつ たばねて かぞえられる。 & \pageref{bk:1-2:tabaneru} & \bkcb & \bkcb \\
 & 23の 2と 3の いみが わかる。 & \pageref{bk:1-2:kurai} & \bkcb & \bkcb \\
 & 100、 1000、 10000が わかる。 & \pageref{bk:1-2:hyaku} & \bkcb & \bkcb \\
{\sffamily\bfseries\color{bkTheme}1-3} & なにも ない ことを 0で あらわせる。 & \pageref{bk:1-3:nai} & \bkcb & \bkcb \\
 & 0に あわせて ながさを よめる。 & \pageref{bk:1-3:awaseru} & \bkcb & \bkcb \\
 & なにも ない くらいに 0を かける。 & \pageref{bk:1-3:kurai} & \bkcb & \bkcb \\
{\sffamily\bfseries\color{bkTheme}1-4} & 「あわせる」「ふえる」を しきに かける。 & \pageref{bk:1-4:awaseru} & \bkcb & \bkcb \\
 & ものさしの うえで たし\R{算}{ざん}が できる。 & \pageref{bk:1-4:monosashi} & \bkcb & \bkcb \\
 & 10を こえる たし\R{算}{ざん}が できる。 & \pageref{bk:1-4:kuriagari} & \bkcb & \bkcb \\
{\sffamily\bfseries\color{bkTheme}1-5} & 「のこり」「ちがい」を しきに かける。 & \pageref{bk:1-5:nokori} & \bkcb & \bkcb \\
 & どれだけ かわったかを もとめられる。 & \pageref{bk:1-5:kawatta} & \bkcb & \bkcb \\
 & 10の まとまりを くずして ひける。 & \pageref{bk:1-5:kuzusu} & \bkcb & \bkcb \\
{\sffamily\bfseries\color{bkTheme}1-6} & \R{数直線}{すう|ちょく|せん}で たし\R{算}{ざん}・ひき\R{算}{ざん}が できる。 & \pageref{bk:1-6:susumu} & \bkcb & \bkcb \\
 & 2つの かずの ちがいを みつけられる。 & \pageref{bk:1-6:aida} & \bkcb & \bkcb \\
 & 1\R{目}{め}もりが 10や 100でも よめる。 & \pageref{bk:1-6:ookii} & \bkcb & \bkcb \\
{\sffamily\bfseries\color{bkTheme}1-7} & \R{何十}{なん|じゅう}・\R{何百}{なん|びゃく}の けいさんが できる。 & \pageref{bk:1-7:matomari} & \bkcb & \bkcb \\
 & くらいを そろえて、 ひっ\R{算}{さん}が できる。 & \pageref{bk:1-7:kuraigoto} & \bkcb & \bkcb \\
 & くり\R{上}{あ}がり・くり\R{下}{さ}がりが わかる。 & \pageref{bk:1-7:kuriagari} & \bkcb & \bkcb \\
{\sffamily\bfseries\color{bkTheme}1-8} & 「1\,cm の いくつ\R{分}{ぶん}」で ながさを いえる。 & \pageref{bk:1-8:cm} & \bkcb & \bkcb \\
 & ものさしで cm と mm を よめる。 & \pageref{bk:1-8:yomu} & \bkcb & \bkcb \\
 & ながさの たし\R{算}{ざん}・ひき\R{算}{ざん}が できる。 & \pageref{bk:1-8:keisan} & \bkcb & \bkcb \\
{\sffamily\bfseries\color{bkTheme}1-9} & \R{水}{みず}の かさを、 L・dL・mL で いえる。 & \pageref{bk:1-9:kasa} & \bkcb & \bkcb \\
 & てんびんで、 おもさを くらべられる。 & \pageref{bk:1-9:tenbin} & \bkcb & \bkcb \\
 & おもさを、 g・kg で いえる。 & \pageref{bk:1-9:g} & \bkcb & \bkcb \\
{\sffamily\bfseries\color{bkTheme}1-10} & \mbox{\R{時計}{と|けい}}を よんで、 \R{時}{じ}こくが いえる。 & \pageref{bk:1-10:shikumi} & \bkcb & \bkcb \\
 & \R{時}{じ}こくと \mbox{\R{時間}{じ|かん}}の ちがいが わかる。 & \pageref{bk:1-10:jikan} & \bkcb & \bkcb \\
 & 60で くり\R{上}{あ}がる けいさんが できる。 & \pageref{bk:1-10:rokujuu} & \bkcb & \bkcb \\
{\sffamily\bfseries\color{bkTheme}1-11} & \mbox{1\R{目}{め}もり}が いくつかを \R{見}{み}つけられる。 & \pageref{bk:1-11:hitomemori} & \bkcb & \bkcb \\
 & \mbox{\R{温度計}{おん|ど|けい}}や はかりの \R{目}{め}もりを よめる。 & \pageref{bk:1-11:ondokei} & \bkcb & \bkcb \\
 & \R{目}{め}もりの あいだを 「だいたい」で よめる。 & \pageref{bk:1-11:daitai} & \bkcb & \bkcb \\
{\sffamily\bfseries\color{bkTheme}1-12} & 「1つ\R{分}{ぶん} × いくつ\R{分}{ぶん}」の しきが かける。 & \pageref{bk:1-12:matomari} & \bkcb & \bkcb \\
 & かけ\R{算}{ざん}を おはじきの \R{四角}{し|かく}で あらわせる。 & \pageref{bk:1-12:shikaku} & \bkcb & \bkcb \\
 & 2のだん〜5のだんの \R{九九}{く|く}を つかえる。 & \pageref{bk:1-12:dan} & \bkcb & \bkcb \\
{\sffamily\bfseries\color{bkTheme}1-13} & \R{九九}{く|く}の \R{表}{ひょう}で、 こたえを \R{見}{み}つけられる。 & \pageref{bk:1-13:yomikata} & \bkcb & \bkcb \\
 & \R{表}{ひょう}の きまりを \R{見}{み}つけられる。 & \pageref{bk:1-13:fuekata} & \bkcb & \bkcb \\
 & \R{入}{い}れかえても、 こたえは \R{同}{おな}じと わかる。 & \pageref{bk:1-13:irekae} & \bkcb & \bkcb \\
{\sffamily\bfseries\color{bkTheme}1-14} & \R{分}{わ}ける ことを、 わり\R{算}{ざん}の しきに かける。 & \pageref{bk:1-14:hitoribun} & \bkcb & \bkcb \\
 & わり\R{算}{ざん}の こたえを、 \R{九九}{く|く}で \R{見}{み}つける。 & \pageref{bk:1-14:kuku} & \bkcb & \bkcb \\
 & 0で わる わり\R{算}{ざん}は できないと わかる。 & \pageref{bk:1-14:zero} & \bkcb & \bkcb \\
{\sffamily\bfseries\color{bkTheme}1-15} & あまりの ある わり\R{算}{ざん}が できる。 & \pageref{bk:1-15:amari} & \bkcb & \bkcb \\
 & わり\R{算}{ざん}の こたえを たしかめられる。 & \pageref{bk:1-15:tashikame} & \bkcb & \bkcb \\
 & くり\R{返}{かえ}す ものを わり\R{算}{ざん}で しらべられる。 & \pageref{bk:1-15:shuu} & \bkcb & \bkcb \\
{\sffamily\bfseries\color{bkTheme}1-16} & はんぱを \mbox{\R{小数}{しょう|すう}}で あらわせる。 & \pageref{bk:1-16:shousuu} & \bkcb & \bkcb \\
 & \mbox{\R{数直線}{すう|ちょく|せん}}で \mbox{\R{小数}{しょう|すう}}が よめる。 & \pageref{bk:1-16:suuchokusen} & \bkcb & \bkcb \\
 & \mbox{\R{半分}{はん|ぶん}}を \mbox{\R{分数}{ぶん|すう}}で あらわせる。 & \pageref{bk:1-16:bunsuu} & \bkcb & \bkcb \\
{\sffamily\bfseries\color{bkTheme}1-17} & しらべて かぞえ、 \R{表}{ひょう}に まとめられる。 & \pageref{bk:1-17:hyou} & \bkcb & \bkcb \\
 & ぼうグラフが かけて、 よめる。 & \pageref{bk:1-17:bou} & \bkcb & \bkcb \\
 & グラフの 1\R{目}{め}もりを たしかめられる。 & \pageref{bk:1-17:memori} & \bkcb & \bkcb \\
{\sffamily\bfseries\color{bkTheme}1-18} & \R{表}{ひょう}を、 おれ\R{線}{せん}グラフに かける。 & \pageref{bk:1-18:kaku} & \bkcb & \bkcb \\
 & \R{線}{せん}の むきで、 ふえる・へるが わかる。 & \pageref{bk:1-18:muki} & \bkcb & \bkcb \\
 & かたむきで、 かわる はやさが わかる。 & \pageref{bk:1-18:kyuu} & \bkcb & \bkcb \\
{\sffamily\bfseries\color{bkTheme}1-19} & \R{三角形}{さん|かく|けい}と \R{四角形}{し|かく|けい}を \R{見}{み}わけられる。 & \pageref{bk:1-19:sankaku} & \bkcb & \bkcb \\
 & \R{直角}{ちょっ|かく}や \R{長方形}{ちょう|ほう|けい}・\R{正方形}{せい|ほう|けい}が わかる。 & \pageref{bk:1-19:chokkaku} & \bkcb & \bkcb \\
 & \R{円}{えん}の \R{中心}{ちゅう|しん}と \R{半径}{はん|けい}が わかる。 & \pageref{bk:1-19:en} & \bkcb & \bkcb \\
{\sffamily\bfseries\color{bkTheme}1-20} & \R{角}{かく}は 「まわった \R{大}{おお}きさ」と わかる。 & \pageref{bk:1-20:kaku} & \bkcb & \bkcb \\
 & 1\R{回転}{かい|てん}は 360°、 \R{直角}{ちょっ|かく}は 90°と わかる。 & \pageref{bk:1-20:do} & \bkcb & \bkcb \\
 & \R{分度器}{ぶん|ど|き}で \R{角}{かく}を はかれる。 & \pageref{bk:1-20:bundoki} & \bkcb & \bkcb \\
{\sffamily\bfseries\color{bkTheme}1-21} & \R{広}{ひろ}さを、 マスの \R{数}{かず}で くらべられる。 & \pageref{bk:1-21:masu} & \bkcb & \bkcb \\
 & 1\,cm\textsuperscript{2} を もとに、 \R{面積}{めん|せき}を あらわせる。 & \pageref{bk:1-21:cm} & \bkcb & \bkcb \\
 & \R{長方形}{ちょう|ほう|けい}の \R{面積}{めん|せき}を もとめられる。 & \pageref{bk:1-21:kakezan} & \bkcb & \bkcb \\
\end{longtblr}
\endgroup


%% ---------------- つぎのステージへ（第2巻の予告・第5回で追加） ----------------
\bkbacktitle{つぎのステージへ}
この \R{本}{ほん}で、 かぞえる・ はかる・ くらべる ことが できるように なりました。 つぎの 2\R{巻}{かん} 『\R{数}{かず}と \R{量}{りょう}の ことば』では、 こんな ふしぎを しらべます。\par\vspace{3mm}
\newcommand\bkqzcard[4]{% 章, 色, 問い, 絵
  \begin{tikzpicture}[x=1mm, y=1mm]
    \fill[white, draw=bkVol2, line width=1pt, rounded corners=3mm] (0,0) rectangle (66,48);
    \node[fill=bkVol2, text=bkOn2, rounded corners=1mm, font=\sffamily\bfseries\footnotesize, inner xsep=1.6mm, inner ysep=0.8mm] at (9,43) {#1};
    \begin{scope}[shift={(33,24)}] #4 \end{scope}
    \node[anchor=south, text width=60mm, align=center, font=\fontsize{12}{19}\selectfont, text=bkInk] at (33,2.5) {#3};
  \end{tikzpicture}}
\begin{bkcenter}
\bkqzcard{2-2}{}{0より \R{小}{ちい}さい かず って ある？}{%
  \draw[bkInk, line width=0.8pt] (-26,2) -- (26,2);
  \foreach \k in {-5,...,5} {\draw[bkInk, line width=0.6pt] (\k*4.6,0.8) -- (\k*4.6,3.2);}
  \node[font=\sffamily\bfseries\scriptsize] at (0,-1.5) {0};
  \node[font=\sffamily\bfseries\scriptsize, text=bkBlue] at (-13.8,-1.5) {？};
  \node[font=\sffamily\bfseries\scriptsize] at (13.8,-1.5) {3};
  \draw[bkBlue, line width=1pt, -{Stealth[length=2mm]}] (0,6) to[bend right=35] (-13.8,6);}\hfill
\bkqzcard{2-4}{}{ピザを 3\R{人}{にん}で \R{分}{わ}けると？}{%
  \bkpizza{(0,2)}{11}{3}{1}{bkOrange}}\par\vspace{4mm}
\bkqzcard{2-5}{}{「20\%びき」って どれだけ？}{%
  \fill[bkYellow!35, draw=bkStem, line width=0.8pt, rounded corners=1mm] (-18,-6) rectangle (18,10);
  \node[font=\sffamily\bfseries\large, text=bkRed] at (0,4) {20\%};
  \node[font=\sffamily\bfseries\footnotesize, text=bkInk] at (0,-2.5) {わりびき};}\hfill
\bkqzcard{2-11}{}{かずの かわりに \R{文字}{も|じ}？}{%
  \node[font=\fontsize{20}{20}\selectfont, text=bkInk] at (0,2) {□ ＋ 3 ＝ 8};}
\end{bkcenter}
\vspace{3mm}
\bkfukidashi{delta}{1\R{巻}{かん}で のった 7\R{本}{ほん}の \R{路線}{ろ|せん}は、 2\R{巻}{かん}へ つづいて いるよ。 つぎの ステージで まってるね！}

%% ---------------- ふろく（コピーして つかう ページ・第5回で追加） ----------------
%% 見開き（p.236〜237）。 ものさしと 方眼は 実寸（x=1mm）。
\tikzset{bk cut/.style={draw=bkGray, line width=0.5pt, dash pattern=on 1.6pt off 1.3pt}}
\bkbacktitle{ふろく}
{\small コピーして つかう ページです。 コピーは 「100％」で とろう。 \R{点線}{てん|せん}は きる ところ。 はさみは、 おとなと いっしょに つかおう。}\par

\bkuhead{10の わくと おはじき}{（1-2・1-4）}
\begin{bkcenter}
\begin{tikzpicture}[x=1mm, y=1mm]
  \foreach \f in {0,1} {\begin{scope}[shift={(0,-\f*34)}]
    \fill[white] (0,0) rectangle (70,-28);
    \draw[bkInk, line width=0.6pt] (0,-14) -- (70,-14);
    \foreach \x in {14,28,42,56} {\draw[bkInk, line width=0.6pt] (\x,0) -- (\x,-28);}
    \draw[bkInk, line width=1.2pt] (0,0) rectangle (70,-28);
  \end{scope}}
  \begin{scope}[shift={(78,0)}]
    \foreach \i in {0,...,4} {\foreach \j in {0,...,3} {
      \ifnum\j<2 \colorlet{bk@oc}{bkAppleRed}\else\colorlet{bk@oc}{bkBlue}\fi
      \filldraw[fill=bk@oc!30, draw=bk@oc!80!black, line width=0.5pt] (\i*12.4+6.2,-\j*15.5-7.75) circle (4.6);
      \draw[bk cut] (\i*12.4+6.2,-\j*15.5-7.75) circle (5.9);}}
  \end{scope}
\end{tikzpicture}
\end{bkcenter}
{\small \R{上}{うえ}の \R{段}{だん}が いっぱいで 5、 2\R{段}{だん} いっぱいで 10。 おはじきを おいて、 たし\R{算}{ざん}・ひき\R{算}{ざん}を して みよう。}\par

\bkuhead{かずカード}{（1-1〜1-3）}
\begin{bkcenter}
\begin{tikzpicture}[x=1mm, y=1mm]
  \foreach \k in {0,...,11} {\pgfmathtruncatemacro\cx{mod(\k,6)}\pgfmathtruncatemacro\cy{div(\k,6)}
    \begin{scope}[shift={(\cx*23.8,-\cy*29.5)}]
      \fill[white, rounded corners=1.5mm] (0,0) rectangle (21,-27);
      \draw[bk cut, rounded corners=1.5mm] (0,0) rectangle (21,-27);
      \ifnum\k<11
        \node[font=\fontsize{30}{30}\selectfont, text=bkInk] at (10.5,-10.5) {\k};
        \foreach \d in {0,...,9} {\pgfmathtruncatemacro\dx{mod(\d,5)}\pgfmathtruncatemacro\dy{div(\d,5)}
          \ifnum\d<\k \fill[bkVol] (3.1+\dx*3.7,-20.4-\dy*3.4) circle (1.2);
          \else \draw[bkGray!60, line width=0.4pt] (3.1+\dx*3.7,-20.4-\dy*3.4) circle (1.2);\fi}
      \else
        \node[font=\footnotesize, text=bkGray, align=center] at (10.5,-14) {じぶんで\\かこう};
      \fi
    \end{scope}}
\end{tikzpicture}
\end{bkcenter}

\clearpage
\bkuhead{ものさし}{（1-8）}
\begin{bkcenter}
\begin{tikzpicture}[x=1mm, y=1mm]
  \fill[bkYellow!22, rounded corners=1.5mm] (-5,0) rectangle (127,-19);
  \draw[bk cut, rounded corners=1.5mm] (-5,0) rectangle (127,-19);
  \foreach \m in {0,...,120} {\pgfmathsetmacro\bk@tl{ifthenelse(mod(\m,10)==0,6.5,ifthenelse(mod(\m,5)==0,4.6,3))}
    \draw[bkInk, line width=0.35pt] (\m,0) -- (\m,-\bk@tl);}
  \foreach \c in {0,...,12} {\node[font=\sffamily\bfseries\small, text=bkInk] at (\c*10,-10.6) {\c};}
  \node[font=\sffamily\footnotesize, text=bkInk, anchor=east] at (125,-16) {cm};
\end{tikzpicture}
\end{bkcenter}
{\small コピーしたら、 ほんとうの ものさしと ならべて、 1\,cm が おなじか たしかめよう。}\par

\bkuhead{とけい}{（1-10）}
\noindent\begin{minipage}[c]{0.56\linewidth}\centering
\begin{tikzpicture}[x=1mm, y=1mm]
  \fill[white] (0,0) circle (36.5);
  \draw[bk cut] (0,0) circle (36.5);
  \draw[bkInk, line width=1.2pt] (0,0) circle (30);
  \foreach \m in {0,...,59} {\pgfmathsetmacro\bk@a{90-6*\m}
    \ifnum\numexpr\m-5*(\m/5)\relax=0 \draw[bkInk, line width=1pt] (\bk@a:26.6) -- (\bk@a:30);
    \else \draw[bkInk, line width=0.4pt] (\bk@a:28.4) -- (\bk@a:30);\fi}
  \foreach \h in {1,...,12} {\node[font=\sffamily\bfseries\fontsize{13}{13}\selectfont, text=bkInk] at ({90-30*\h}:21.5) {\h};}
  \foreach \h in {1,...,12} {\pgfmathtruncatemacro\bk@mv{5*\h}
    \node[font=\sffamily\bfseries\fontsize{8.5}{9}\selectfont, text=bkBlue!85!black, inner sep=0pt] at ({90-30*\h}:33.4) {\bk@mv};}
  \fill[bkInk] (0,0) circle (1.1);
\end{tikzpicture}
\end{minipage}\hfill
\begin{minipage}[c]{0.42\linewidth}
\begin{tikzpicture}[x=1mm, y=1mm]
  % ながい はり（ふん）
  \begin{scope}[shift={(4,0)}]
    \draw[bk cut, rounded corners=1mm] (-4,-4.4) rectangle (31,4.4);
    \filldraw[fill=bkRed!25, draw=bkRed!85!black, line width=0.8pt, line join=round] (-2.4,-2.2) -- (22,-1.6) -- (27,0) -- (22,1.6) -- (-2.4,2.2) -- cycle;
    \draw[bkRed!85!black, line width=0.6pt] (0,0) circle (0.9);
    \node[anchor=north west, font=\footnotesize, text=bkInk, inner sep=0pt] at (-4,-6) {ながい はり（\R{分}{ふん}）};
  \end{scope}
  % みじかい はり（じ）
  \begin{scope}[shift={(4,-22)}]
    \draw[bk cut, rounded corners=1mm] (-4,-4.8) rectangle (22,4.8);
    \filldraw[fill=bkInk!30, draw=bkInk, line width=0.8pt, line join=round] (-2.4,-2.6) -- (13.5,-2.2) -- (18,0) -- (13.5,2.2) -- (-2.4,2.6) -- cycle;
    \draw[bkInk, line width=0.6pt] (0,0) circle (0.9);
    \node[anchor=north west, font=\footnotesize, text=bkInk, inner sep=0pt] at (-4,-6.4) {みじかい はり（\R{時}{じ}）};
  \end{scope}
\end{tikzpicture}\par\vspace{2mm}
{\small はりを きりとって、 ○の ところに あなを あけ、 とけいの まん\R{中}{なか}に わりピンで とめよう。}
\end{minipage}

\bkuhead{1\,cm の \R{方眼}{ほう|がん}}{（1-19・1-21）}
\begin{bkcenter}
\begin{tikzpicture}[x=1mm, y=1mm]
  \fill[white] (0,0) rectangle (130,-40);
  \draw[bkGray!70, line width=0.4pt] (0,-40) grid[step=10] (130,0);
  \draw[bkInk, line width=0.8pt] (0,0) rectangle (130,-40);
\end{tikzpicture}
\end{bkcenter}
{\small 1マスが 1\,cm\textsuperscript{2}。 \R{形}{かたち}を かいて、 マスを かぞえて \R{広}{ひろ}さを くらべよう。}

%% ---------------- さくいん ----------------
\bkbacktitle{さくいん}
{\small ことばの あとの \R{数}{かず}は、 ページです。}\par\vspace{2mm}
\begin{multicols}{2}\fontsize{12}{20.5}\selectfont\setlength{\parskip}{0pt}\ltjsetruby{stretch=001}% ルビが長い語も、親字を左にそろえる
\directlua{bk.printidx()}
\end{multicols}

%% ---------------- 奥付 ----------------
\bkpadfour
\global\bk@decofalse\bkfolio{false}
\null\vfill
\noindent\begin{minipage}{\linewidth}\fontsize{10}{17}\selectfont
{\sffamily\bfseries\Large\color{bkDeep}\R{物理}{ぶつ|り}を \R{理解}{り|かい}するための \R{数学}{すう|がく}　\R{第}{だい}1\R{巻}{かん}}\par\vspace{1mm}
{\sffamily\bfseries\large かぞえる・はかる}\par\vspace{3mm}
{\color{bkVol}\hrule height 1pt}\vspace{3mm}
\begin{tabular}{@{}l@{\hspace{6mm}}l@{}}
\R{編者}{へん|しゃ} & 「\R{物理}{ぶつ|り}を \R{理解}{り|かい}するための \R{数学}{すう|がく}」\R{編集委員会}{へん|しゅう|い|いん|かい} \\
 & \R{委員長}{い|いん|ちょう}　\R{小学校長}{しょう|がっ|こう|ちょう}（まとめ\R{役}{やく}） \\
 & \R{委員}{い|いん}　\R{数学}{すう|がく}・\R{電気工学}{でん|き|こう|がく}・\R{機械工学}{き|かい|こう|がく}・\R{教育学}{きょう|いく|がく}の \R{大学教授}{だい|がく|きょう|じゅ}、\\
 & \phantom{\R{委員}{い|いん}　}\R{幼稚園教諭}{よう|ち|えん|きょう|ゆ}、 \R{特別支援学校教諭}{とく|べつ|し|えん|がっ|こう|きょう|ゆ}、 \R{心理士}{しん|り|し}、\\
 & \phantom{\R{委員}{い|いん}　}\R{科学解説者}{か|がく|かい|せつ|しゃ}、 \R{小学校教諭}{しょう|がっ|こう|きょう|ゆ} \\
デザイン & ユニバーサルデザイン\R{担当委員}{たん|とう|い|いん} \\
\R{挿絵}{さし|え} & イラスト\R{担当委員}{たん|とう|い|いん} \\
\R{版}{はん} & \R{見本版}{み|ほん|ばん}（\R{第}{だい}5\R{回編集委員会}{かい|へん|しゅう|い|いん|かい}）　2026\R{年}{ねん}10\R{月}{がつ} \\
\end{tabular}\par\vspace{3mm}
{\color{bkVol}\hrule height 1pt}\vspace{2mm}
この \R{本}{ほん}の もんだいの こたえは、 うしろの 「こたえ」に あります。 くわしい とき\R{方}{かた}は、 べっさつの もんだいしゅうと、 \R{先生}{せん|せい}むけの \R{本}{ほん}に のせます。
\end{minipage}
\makeatother

\endgroup
